% use \chaptermark{}
% to alter or adjust the chapter heading in the running head



\section{Introduction to Semigroup Theory}
\label{sec:2.1}
% BEGIN SUBJECT INDEX: sec:2.1
\index{semigroup!definition}
\index{semigroup!examples}
% END SUBJECT INDEX


For our initial purposes, the definition of a semigroup will be
very straightforward.  No topological notions will be included.

\begin{definition}   \label{defn:2.1}
 A set of elements $S$ equipped with a binary operation
     $*$ will be called a semigroup if
\begin{enumerate}
\item[(a)] for all $s,t \in S$, $s * t \in S$ (closure);

\item[(b)] for all $s,t,u \in S$, $(s * t) * u = s * (t*u)$
     (associativity);
\end{enumerate}
\end{definition}

\noindent
(we will usually drop the $*$ notation and write $s*t = st$).

We will be utilizing certain canonical examples.

\begin{example}	\label{exm:2.1}
Let $X$ be a set and let $\F_X$ be the set of all functions on the set $X$ with the operation composition.  
We will write composition in the following way
\[
  f,g \in \F_X.
\]
Then,
\[
  (x)f \circ g = ((x)\F)g.
\]
Note this is opposite the standard direction for $\F \circ g$ and is done due to the matrix representation of $\F_X$, when $X$ is finite.
(See Example \ref{exm:2.2}.)
\end{example}

\runinhead{Note.}  $X = \{1,2,3\}$
\[
  \begin{matrix}
         &&f           &&   &&g           &&  \\
      1  &&\rightarrow &&1  &&\rightarrow &&1 \\
      2  &&\searrow    &&2  &&\rightarrow &&2 \\
      3  &&\rightarrow &&3  &&\nearrow    &&3 \\
     \end{matrix} \qquad \qquad \qquad
     \begin{matrix}
         &&f \circ g   &&  \\
      1  &&\rightarrow &&1 \\
      2  &&\rightarrow &&2 \\
      3  &&\nearrow    &&3 \\
     \end{matrix}
\]

     An extremely useful concept in all of algebra is that of a
representation.  
Due to the usefulness of matrices, it is often preferable to find an algebraically equivalent matrix semigroup to work with.

\begin{definition}   \label{defn:2.2}
Let $S$ be a semigroup.  
Then a semigroup of matrices $M$ is called a representation of $S$, if there exists a homomorphism $T: S \to M$.  
If $T$ is an isomorphism, then the representation $M$ is called \textsl{faithful}.
\end{definition}

\begin{example}	\label{exm:2.2}
Let $X = \{1,2,\ldots,n\}$ and let $S = \F_X$.

For $f \in \F_X$ define
\[
  [T(f)]_{ij} = \begin{cases}
                    1  &(i)f = j \\
                    0  &\text{otherwise.}
                    \end{cases}
\]
Let $M = \{T(f) \,|\, f \in \F_X\}$.
\end{example}

\runinhead{Homework.}  
Show that $T: S \to M$ is a semigroup
isomorphism.  That is, show that $M$ is a faithful representation of
$S$.

Let $X = \{1,2\}$.
\[
\begin{aligned}
     &\hbox{card}(\F_X) = 4                       \\
%    \vspace{2\jot}
     &\F_1 \ \begin{matrix}
             1 &\rightarrow &1 \\
             2 &\rightarrow &2
             \end{matrix} \qquad \qquad
      \F_2 \ \begin{matrix}
             1 &\searrow &1 \\
             2 &\nearrow &2
             \end{matrix} \qquad \qquad
      \F_3 \ \begin{matrix}
             1 &\rightarrow &1 \\
             2 &\nearrow    &2
             \end{matrix} \qquad \qquad
      \F_4 \ \begin{matrix}
             1 &\searrow    &1 \\
             2 &\rightarrow &2
             \end{matrix}                           \\
%    \vspace{2\jot}
     &T(\F_1) = \begin{bmatrix}
                1  &0  \\  0  &1  \end{bmatrix} \qquad
      T(\F_2) = \begin{bmatrix}
                0  &1  \\  1  &0  \end{bmatrix} \qquad
      T(\F_3) = \begin{bmatrix}
                1  &0  \\  1  &0  \end{bmatrix} \qquad
      T(\F_4) = \begin{bmatrix}
                0  &1  \\  0  &1  \end{bmatrix}
     \end{aligned}
\]

\runinhead{Note.}  The set of matrices $M$ in the above example are also binary (i.e., 0,1) stochastic matrices.  
Their row sums are trivially 1.

\runinhead{Some notation.}  
If $s \in S$, and $A,B$ are subsets of $S$ then
\begin{enumerate}
\item[(1)] $sA = \{sa \,|\, a \in A \}$
\item[(2)] $AB = \{ab \,|\, a \in A, \ a \in B\}$
\end{enumerate}

\begin{definition}   \label{defn:2.3}
\begin{enumerate}
\item[(1)] $T \subset S$ is a subsemigroup of $S$ if
     $T \cdot T \subset T$.  If $T$ is a group, we will call
     $T$ a subsgroup of $S$.

\item[(2)] A subset $L$ of $S$ is called a left ideal, if
     $SL \subset L$.  Right ideals are similarly defined
     ($RS \subset R$).  Note that $L,R$ are necessarily subsemigroups
     of $S$.

\item[(3)] A subset $I$ of $S$ is a (two-sided) ideal if it is both a
     right ideal and a left ideal.

\item[(4)] $S$ is said to be left (right) simple if it contains no
     proper left (right) ideals.  If $S$ contains no proper ideal,
     then $S$ is said to be simple.  ($S$ is trivially an ideal.)
\end{enumerate}
\end{definition}

\runinhead{Homework.}  If $I_1$, $I_2$ are ideals, show that
     $I_1 \cap I_2$ is an ideal.  (In particular, it is not empty.)

\runinhead{Homework.}  Show that $S$ is left simple iff for any
     $a,b \in S$ there exists $s \in S$ such that $sa = b$.
     (Hint: Note that for any $a \in S$, $Sa$ is a left ideal of $S$.)

\begin{enumerate}
\item[(5)] A left ideal is called principal if it is of the form
     $\{a\} \cup Sa$ for some $a \in S$.

\item[(6)] An element $u \in S$ is idempotent if $a^2 = a$.  If an
     idempotent $e$ is such that $es = s$ ($se = s$) then $e$ is
     called a left (right) identity.  If $e$ is both a left and right
     identity, then it is called the identity of $S$.  It is easy to 
     see that an identity must be unique.

\item[(7)] An element $z$ is called a left (right) zero if $zs = z$
     ($sz = z$) for all $s \in S$.  $z$ is a zero, if it is both a left
     and right zero.
\end{enumerate}

\runinhead{Homework.}  Show that a semigroup $S$ is a group iff it is
     both left simple and right simple.

\begin{example}	\label{exm:2.3}
Let $S = \left\{ \begin{bmatrix} a &b \\ b &a
     \end{bmatrix} \, \big| \, a,b \in GF(q), \ a \not= b, \ q = p^n
     \right\}$.  Then $S$ is an abelian semigroup.  (Show this.)
\[
     S = \left\{
         \begin{bmatrix} 1 &0 \\ 0 &1 \end{bmatrix}, \
         \begin{bmatrix} 2 &0 \\ 0 &2 \end{bmatrix}, \
         \begin{bmatrix} 0 &1 \\ 1 &0 \end{bmatrix}, \
         \begin{bmatrix} 0 &2 \\ 2 &0 \end{bmatrix}, \
         \begin{bmatrix} 2 &1 \\ 1 &2 \end{bmatrix}, \
         \begin{bmatrix} 1 &2 \\ 2 &1 \end{bmatrix} \right\}
\]
Notice
\[
\begin{aligned}
     &\det \begin{bmatrix} 2 &1 \\ 1 &2 \end{bmatrix}
            = 3 \equiv 0 \hbox{ mod } 3               \\
     &\det \begin{bmatrix} 1 &2 \\ 2 &1 \end{bmatrix}
            = 0 \hbox{ mod } 3                        \\
     \end{aligned}
\]
In fact, both matrices are rank 1, mod 3.  
Notice also, that $G = \left\{ \begin{bmatrix} 2 &1 \\ 1 &2 \end{bmatrix}, \
\begin{bmatrix} 1 &2 \\ 2 &1 \end{bmatrix} \right\}$ forms an ideal in $S$, has
no ideals within it ($G$ is a minimal ideal) and in fact, $G$ is a group.
\end{example}

\begin{example}		\label{exm:2.4}  
{\bf This example will be of paramount importance!!}

Let $X,Y$ be sets and $G$ a group with $\phi: Y \times X \to G$ any
function.  Let $S = X \times G \times Y = \{(x,g,y)\}$ define
\[
  (x_1,g_1,y_1)(x_2,g_2,y_2) = (x_1,g_1 \phi(y_1,x_2) g_2,y_2).
\]
Then $S$ becomes a semigroup under this product.
\end{example}

We will begin our next lecture by discussing this example at length.
